% Inserted after the trajectory section.
% Chunks that read: composition of contributions, and what it costs.

\section{Contributions that compose}
\label{sec:composition}

As presented so far a chunk is a nullary closure: it is evaluated and emits,
and what it emits cannot depend on what any other node produced. Every
contribution therefore stands alone, and a node's values are a collection of
independent results that happen to share an identity.

This section widens a chunk to declare a \emph{read set}. A chunk may name
identities whose values it consults, so that a node's value can be a function
of other nodes' values. The change is what allows several lines of work to
converge into one conclusion rather than merely to accumulate at one address.
It costs three things, and we take them in turn.

\begin{definition}[Read set; read graph]
\label{def:reads}
A chunk carries a finite set of identities $\mathrm{reads}(c)\subseteq\Tau$,
declared at contribution. The \emph{read graph} of a protocol has an edge
$\tau\to\tau'$ whenever some chunk on the node $\tau$ declares a read of
$\tau'$.
\end{definition}

The read set is declared rather than discovered. The kernel requires it before
executing anything --- to order execution, to refuse a contribution that would
close a cycle, and to record which reads were permitted --- and it cannot
obtain it by inspecting a closure.

\begin{definition}[Restricted view]
\label{def:view}
When a chunk is evaluated it receives a \emph{view}: a map from each declared
identity to that node's value collection. Identities absent from the read set
are absent from the view, whatever the graph contains.
\end{definition}

\begin{proposition}[The declared read set bounds what a chunk observes]
\label{prop:read-bound}
A chunk's emitted value is a function of its declared reads and of sources
outside the graph; it is not a function of any node it did not declare.
Consequently the read graph is a sound over-approximation of the run's
read-to-emit relation.
\end{proposition}

\begin{proof}
By \Cref{def:view} the view supplied to a chunk contains exactly the declared
identities, so no undeclared node's values are available to it. The emitted
value is therefore determined by the view together with whatever external
source the closure consults, and every edge realised in the run appears in the
read graph.
\end{proof}

\Cref{prop:read-bound} recovers, in part, something \Cref{thm:emergence}
denied. The trajectory remains not a function of the protocol --- external
sources still vary between runs, and a chunk may consult a declared node
without its value affecting the outcome --- but the read graph now bounds
which trajectories are possible. What was formerly an obligation on modules to
log their reads is now structure the protocol carries.

\subsection{Cycles are refused}

\begin{definition}[Runnable]
\label{def:runnable}
A pending chunk is \emph{runnable} when every identity it declares reading has
produced at least one value. A node is ready when it carries a runnable
pending chunk.
\end{definition}

Readiness is thereby data-dependent, where before it was a property of the
chunk bag alone. A protocol whose read graph contains a cycle has no
execution order: each node in the cycle waits on another, and none becomes
runnable.

\begin{theorem}[Acyclicity is necessary and sufficient for progress]
\label{thm:acyclic}
Let a protocol have a finite read graph. Every contributed chunk eventually
becomes runnable if and only if the read graph is acyclic and every declared
identity is contributed.
\end{theorem}

\begin{proof}
($\Leftarrow$) A finite acyclic digraph admits a topological order. Taking
nodes in that order, every identity a chunk reads precedes it and has
therefore produced values by the time the chunk is considered, so the chunk is
runnable.

($\Rightarrow$) Suppose the read graph contains a cycle
$\tau_1\to\tau_2\to\cdots\to\tau_k\to\tau_1$. For any chunk on $\tau_1$
participating in the cycle to be runnable, $\tau_2$ must have produced values;
for that, $\tau_3$ must have; and so on around the cycle to $\tau_1$ itself,
which by hypothesis has not. No member of the cycle is ever runnable. If
instead some declared identity is never contributed, the chunks reading it are
never runnable by \Cref{def:runnable}.
\end{proof}

The kernel therefore refuses, at contribution time, any contribution that
would close a cycle in the read graph, and reports the cycle it found. Three
remarks on why refusal is the right response rather than the alternatives.

\begin{remark}[Refusal is structural, not adjudication]
\label{rem:refusal-inert}
The check inspects declared read sets and never values, so
\Cref{thm:inertia} is untouched. The kernel is not judging a contribution to
be wrong; it is reporting that no execution order exists for it. A refused
contribution leaves the graph unchanged, so a contributor may correct the
declaration and retry.
\end{remark}

\begin{remark}[Why not a bounded fixed-point iteration]
\label{rem:no-fixpoint}
A cyclic read graph could be admitted and iterated to a fixed point under a
step or resource bound. We decline this because it would make the terminal
values a function of that bound, so two runs of one protocol under different
bounds would disagree --- and the separation between what can be computed and
what was (\Cref{cor:repro}) rests on the protocol determining the former.
A construct whose result depends on how much resource it was given belongs to
the trajectory, not the protocol, and should not masquerade as the latter.
\end{remark}

\begin{remark}[Acyclicity does not conflict with mutual support]
\label{rem:not-coherence}
The coherence requirement of the companion language --- that a claim rest on
at least three \emph{mutually independent} sources --- might appear to demand
a cycle. It does not, and the distinction is worth stating because the
opposite reading is natural.

That requirement concerns a \emph{support} graph: an observer's assessment,
after the fact, of which findings corroborate which. Mutual independence
means precisely that the sources do \emph{not} inform one another. Three
catalysts reading each other's outputs would violate the requirement rather
than satisfy it. The read graph and the support graph are different objects,
and acyclicity of the first is consistent with --- indeed enforces --- the
independence the second requires.
\end{remark}

\begin{proposition}[Unreachable reads are reported, not silently stalled]
\label{prop:blocked}
A pending chunk whose declared reads are unsatisfied is reported together with
the identities it awaits, and with whether any of them was never contributed.
A sweep containing only such chunks performs no executions and terminates.
\end{proposition}

The distinction the report carries is between waiting on a node that exists
but has not run, which a further sweep may resolve, and waiting on an identity
nobody contributed, which cannot resolve without a further contribution.
Neither is an error the kernel adjudicates; both are conditions a module may
wish to act on.

\subsection{A node's values remain a collection}

Several chunks may emit onto one node, and a consumer reading that node
receives all of their values. It does not receive a single answer, and the
kernel provides no means of obtaining one.

\begin{remark}[Why no reduction]
\label{rem:no-reduction}
Reducing a collection to one value requires a criterion --- the mean, the
latest, the majority, the most confident --- and every such criterion is a
judgement about which values count. Supplying one in the kernel would place
that judgement in the substrate, where it is invisible and applies uniformly
to payloads it was never chosen for.

The consequence for a consumer is real: every reader must handle a
collection. We take that cost deliberately. A node whose values agree and a
node whose values conflict are different situations, and an interface that
returns one number for both has discarded the difference at the point where
it was still recoverable.
\end{remark}

The two situations are the two an inquiry already admits. A collection whose
values agree, under a criterion the consumer names, is a convergent closure; a
collection whose values disagree is a contested closure, whose honest report
is the plurality together with what produced each member. A module performing
that classification is choosing a criterion visibly, which is where the choice
belongs.

\begin{remark}[A majority is not agreement]
\label{rem:majority}
It is tempting to resolve a contested collection by its modal value when the
mode is large enough. We report the modal share and do not use it to resolve.
The reason is the one that tells against threshold criteria generally: a
majority rule discards the minority precisely in the cases where the minority
is the informative part, and it does so without recording that it did. A collection in which four sources agree and one
dissents is contested, and the dissent is the thing a reader most needs to
see.
\end{remark}
